% GENERATED by scripts/prepare_usenix_source.py; do not edit.
This appendix is written for the USENIX Security 2027 submission requirement. It must be updated after the release-environment run and before anonymous upload. Every statement below distinguishes currently available material from planned release material.

\subsection{A. Artifact availability}

At submission, the authors plan to provide an anonymous repository containing:

\begin{itemize}
\item the MCPGate implementation and integrated filesystem mediator;
\item unit, adversarial, ablation, concurrency, and baseline tests;
\item compact raw JSON sufficient to regenerate each paper table and figure;
\item exact Python and third-party MCP-server dependency locks;
\item the Dockerfile, image identifier, invocation metadata, and build/run commands;
\item annotation codebook, blank forms, anonymized labels, and scoring code; and
\item scripts for quick evaluation and the complete evaluation.
\end{itemize}

The current working repository is not the anonymous artifact. It contains Git history and paths that may identify the author and therefore must not be linked from a double-blind submission.

\subsection{B. Claims supported by the artifact}

The quick artifact supports:

\begin{enumerate}
\item the executable transcript-indistinguishability illustration;
\item effect-oracle separation;
\item controlled response-auditor/gateway baselines;
\item integrated mediator invariants and adversarial regressions;
\item controlled property ablations;
\item local child-process concurrency; and
\item reference and submission-integrity checks.
\end{enumerate}

The full artifact is intended to support:

\begin{enumerate}
\item the 8,692 \ensuremath{\rightarrow} 4,121 \ensuremath{\rightarrow} 1,242 corpus funnel;
\item response-auditor operating points;
\item contract expressibility over the frozen server set;
\item the pinned third-party-server integrated mediator run;
\item held-out baseline and robustness evaluation; and
\item all final performance and confidence-interval figures.
\end{enumerate}

\subsection{C. Quick evaluation}

Expected runtime on the development machine is under two minutes after dependencies are installed:

\begin{verbatim}
python scripts/run_quick_artifact.py
\end{verbatim}

The command creates a timestamped result directory with:

\begin{itemize}
\item exact command lines and exit codes;
\item stdout and stderr logs;
\item wall-clock duration;
\item Git commit and dirty-state indicator;
\item Python/platform and direct-package versions; and
\item a SHA-256 manifest.
\end{itemize}

The checked-in development run is not a release result because its Git state is dirty and it was produced on Windows. The anonymous release must contain a clean, pinned Linux run.

\subsection{D. Full evaluation}

The full evaluation must run on an isolated Linux/Docker host. Unvetted MCP packages must never run directly on a credential-bearing host. The integrated real-server runner:

\begin{itemize}
\item builds a fresh Git-revision/version-tagged image;
\item verifies the locked npm package version and registry SHA-512 integrity;
\item disables runtime networking;
\item uses a read-only container root and bounded tmpfs workspaces;
\item drops all capabilities before adding only those required by the trusted adapter; and
\item records the Git state, Docker version, image ID, host platform, exact package version/integrity, and complete command.
\end{itemize}

The release instructions must state expected hardware, disk use, network use, and wall-clock time separately for quick and full modes.

\subsection{E. Data and provenance}

For each reported table and figure, the release must provide:

\begin{itemize}
\item compact raw rows or a documented aggregate when redistribution is not possible;
\item a schema that defines attempted, protocol-error, authorized-effect, unauthorized-effect, completion, and UNKNOWN fields independently;
\item exact input hashes and byte lengths;
\item the transformation command and software environment; and
\item the regenerated table/figure hash.
\end{itemize}

Existing hashes of unavailable local traces are provenance aids, not a substitute for data release. The paper will not claim fresh-clone reproducibility until an evaluator can regenerate every result named in the claim--evidence matrix.

\subsection{F. Annotation materials}

The release will include the frozen 265-item Round-2 sample, labeling codebook, two independently completed anonymized label files, scorer, agreement result, adjudication policy, and classifier report. If redistribution of source text is restricted, the artifact will include stable record identifiers, hashes, and the maximum redistributable context, with the limitation stated explicitly.

\subsection{G. Licensing and maintenance}

Before release, the authors must audit:

\begin{itemize}
\item the repository license;
\item licenses for copied/adapted scripts and third-party server packages;
\item whether registry metadata and server identifiers may be redistributed; and
\item a durable archival location and DOI for the camera-ready artifact.
\end{itemize}

No archival DOI or long-term maintenance commitment is claimed yet.

\subsection{H. Current reproducibility gaps}

As of 2026-09-25, the following remain open:

\begin{itemize}
\item the integrated Docker result must be repeated from a clean release commit;
\item several historical raw M3--M5 traces are gitignored;
\item the frozen 10-server workloads have not completed the matched five-condition evaluation; and
\item Round-2 human annotation is blank.
\end{itemize}

These are submission blockers, not post-acceptance cleanup items.
